\section{Datos}
\label{sec:datos}

Se eligieron tres tareas que varían en número de clases, longitud de texto y
tamaño, para que las conclusiones no dependan de un único régimen: SST-2
\cite{socher2013sst}, distribuida dentro del \textit{benchmark} GLUE
\cite{wang2019glue}, y AG News y Yelp Polarity, ambas en la recopilación de
Zhang et al.\ \cite{zhang2015character}.

\begin{table}[h]
\centering
\small
\begin{tabular}{llcrrrc}
\toprule
\textbf{Dataset} & \textbf{Fuente (HuggingFace)} & \textbf{Clases} &
\textbf{Train} & \textbf{Val.} & \textbf{Test} & \textbf{\texttt{max\_length}} \\
\midrule
SST-2          & \texttt{nyu-mll/glue}, conf.\ \texttt{sst2} & 2 & \num{60614}  & \num{6735}  & \num{872}   & 64  \\
AG News        & \texttt{fancyzhx/ag\_news}                  & 4 & \num{108000} & \num{12000} & \num{7600}  & 128 \\
Yelp Polarity  & \texttt{fancyzhx/yelp\_polarity}            & 2 & \num{100000}$^{*}$ & \num{20000}$^{*}$ & \num{38000} & 256 \\
\bottomrule
\end{tabular}
\caption{Los tres datasets, ya normalizados por \texttt{src/data\_adapter.py}.
($^{*}$) Yelp va submuestreado; véase más abajo.}
\end{table}

\subsection{Normalización}

Los tres datasets llegan del Hub con esquemas distintos: SST-2 llama
\texttt{sentence} a su columna de texto y arrastra una columna \texttt{idx};
AG News y Yelp usan \texttt{text}; ninguno de los tres trae los mismos splits.
El módulo \texttt{src/data\_adapter.py} los reduce a una forma única --- un
\texttt{DatasetDict} con splits \texttt{train}/\texttt{validation}/\texttt{test}
y solo las columnas \texttt{text} (str) y \texttt{label} (int) --- de modo que
el resto del \textit{pipeline} no sabe con qué dataset trabaja. Añadir una
tarea nueva es añadir una entrada a un registro, no tocar el entrenamiento.

\subsection{Dos decisiones sobre los splits}

Ambas afectan a cómo deben leerse los números, así que conviene dejarlas
explícitas.

\paragraph{El \texttt{test} de SST-2 no se usa.} En GLUE ese split viene sin
etiquetas (\texttt{label = -1}), porque la evaluación oficial es en servidor.
Se usa entonces el \texttt{validation} oficial (\num{872} frases) como
conjunto de test, y se aparta un \SI{10}{\percent} del train como validación
propia. Es el protocolo habitual en la literatura, y es lo que hace que
nuestros números sean comparables con los publicados. El adaptador verifica
explícitamente que ningún split contenga la etiqueta $-1$ y aborta si la
encuentra: es la clase de error que, si pasa inadvertido, produce métricas sin
sentido sin lanzar ninguna excepción.

\paragraph{Yelp va submuestreado.} Se usan \num{100000} de las \num{560000}
reseñas de train y \num{20000} de validación. Entrenar con el conjunto
completo a 256 \textit{tokens} multiplicaría por cinco el costo para mover la
\textit{accuracy} algunas décimas. Lo importante para la validez de la
comparación es que \textbf{el mismo presupuesto se aplica a los dos modelos}:
el submuestreo usa la misma semilla y produce exactamente el mismo
subconjunto para BERT y para DistilBERT.

En los dos datasets que no traen un split de validación etiquetado (AG News y
Yelp), la validación es un \SI{10}{\percent} del train separado con semilla
fija. El conjunto de test nunca participa en ninguna decisión.
